\documentclass[12pt]{article}
\usepackage[utf8]{inputenc}
\usepackage{amsmath,amssymb}
\usepackage{tcolorbox}
\usepackage{xcolor}
\usepackage{geometry}
\usepackage{paracol} % Manejar varias columnas

\geometry{margin=1.5cm}

% Estilos para las cajas
\tcbset{
	1/.style={colback=black!5!white, colframe=black!75!black, title=Funciones de Supervivencia, width=\linewidth},
	calculo/.style={colback=black!5!white, colframe=black!60!black, title=Cálculo, width=\linewidth},
	probabilidad/.style={colback=black!5!white, colframe=black!85!black, title=Probabilidad, width=\linewidth},
	geometria/.style={colback=black!5!white, colframe=black!70!black, title=Geometría, width=\linewidth},
}

\begin{document}
	
	\begin{center}
		{\LARGE \textbf{Formulario Matemáticas Actuariales I}} \\[1em]
		{\small Formulario para el examen parcial}
	\end{center}
	
	\begin{paracol}{2}
		
		% Columna izquierda
		\begin{tcolorbox}[1]
			$$S_x = 1 - F_x$$
		\end{tcolorbox}
		
		\begin{tcolorbox}[probabilidad]
			\textbf{Probabilidad total:}  
			\[
			P(A) = \sum_i P(A \mid B_i)P(B_i)
			\]
			
			\textbf{Teorema de Bayes:}  
			\[
			P(A \mid B) = \frac{P(B \mid A)P(A)}{P(B)}
			\]
		\end{tcolorbox}
		
		\switchcolumn
		
		% Columna derecha
		\begin{tcolorbox}[calculo]
			\textbf{Derivadas básicas:}  
			\[
			\frac{d}{dx}(x^n) = nx^{n-1}
			\]
			\[
			\frac{d}{dx}(\sin x) = \cos x
			\]
			\[
			\frac{d}{dx}(\cos x) = -\sin x
			\]
			
			\textbf{Integrales básicas:}  
			\[
			\int x^n dx = \frac{x^{n+1}}{n+1}+C
			\]
			\[
			\int e^x dx = e^x + C
			\]
		\end{tcolorbox}
		
		\begin{tcolorbox}[geometria]
			\textbf{Área y volumen:}  
			\[
			A_{\text{círculo}} = \pi r^2, \quad 
			V_{\text{esfera}} = \tfrac{4}{3}\pi r^3
			\]
			
			\textbf{Teorema de Pitágoras:}  
			\[
			a^2+b^2=c^2
			\]
		\end{tcolorbox}
		
	\end{paracol}
	
\end{document}
